\cotefile{29}
\scanpage{31}
% archive sheet number: 30

passage au complété (par exemple si $\underline{O}_{S,s}$ à fibres formelles géom. normales, par exemple s'il est excellent). Les mêmes observations s\unsure{\deleted{d}'a|\deleted{'a}d}ppliqueront sans hypothèse de normalité ou de réduction quelconque; à condition de se borner aux revêtements \deleted{kummériens}\inserted{mod.ram.} \uline{galo\unsure{\deleted{s}i|\deleted{i}s|\deleted{n}i|\deleted{i}n}\unsure{\deleted{d}s|\deleted{s}d|\deleted{i}s|\deleted{s}i}iens} ; cf $n^{\underline{o}}$ suivant.

\unsure{\deleted{3}4|\deleted{4}3}. \uline{Cas des revêtements ga\unsure{\deleted{h}l|\deleted{l}h}oisiens}.

Un revêtement \uline{galoisien} de type $\underline{R}$ de $S$, à groupe de Galois $G$, est un objet $X$ de $\operatorname{Rev}^{\underline{R}}(S)$, sur lequel le groupe fini $G$ opère à \deleted{\unsure{r|n|x}}droite, de tel façon que dans la catégorie $\operatorname{Rev}^{\underline{R}}(S)$ on obtienne un\deleted{e} revêtement principal de l'objet final, i.e. $1^{\underline{o}}$) le morphisme $XxG\quad XxX$ est un isomorphisme et $2^{\underline{o}}$) le morphisme $X\quad\text{objet final}$ est couvrant. Comme $\operatorname{Rev}^{\underline{R}}(S)$ est en tous cas une catégorie multigaloisienne, la notion de revêtement galoisi\unsure{\deleted{i}e|\deleted{e}i|\deleted{t}e|\deleted{e}t}n de type $\underline{R}$ s'explicite également en termes des groupesfondamentaux décrivant cette catégorie. On peut enfin l'expliciter en revenant à la définition de $\operatorname{Rev}^{\underline{R}}(S)$, et en exigeant que les revêtements des $\underline{s}(M)$ $(M\quad\operatorname{Ob}\underline{C}_{S'}\ )$ correspondants à $X$ soient principaux de groupe $G$. Par abus de langage, l\unsure{\deleted{'}e|\deleted{e}'}\deleted{objet}\inserted{revêtement} à opérateurs $(\underline{r}(X),G)$ de $S$ s'appellera également revêtement galoisien de type $\underline{R}$ de $S$. Cet abus de langage ne présente pas d'inconvénient trop sérieux dans le cas où le foncteur de la catégorie des revêtements galoisiens de type $\underline{R}$ de $S$, dans la catégorie des revêtements à opérateurs de $S$, défini par le foncteur ``réalisation géométrique'', est pleinement fidèle,lorsqu'on se borne à prendre comme morphismes dans la première les seuls isomor\deleted{ph}phismes. (Alors, en effet, pour un revêtement à opérateurs $(T,G)$ donné de $S$, il existe, à isomorphisme canonique près, au plus un $(X,G)$ \deleted{dont} galoisien de type $\underline{R}$ dont il provient). Nous allons donner des conditi\deleted{o}ons pour qu'il en soit bien ainsi, en nous plaçant comme d'habitude dans le cas de l'exemple 3. Voici d'abord la condition néc. et suffisan
